\section{Method}
\label{sec:method}

\method has three parts that share one piece of machinery. A \emph{graph traversal}
decides which documents share a training sequence and in what order; a
\emph{link-aware attention mask} decides which tokens may attend across document
boundaries; and a \emph{link-following generation loop} runs the same link detection
and the same mask at inference, fetching a linked document into the context when the
model emits a link to it. The link detector, the target-matching key, the grant
geometry, and the ordering constraint are a single implementation used in both
regimes (\Cref{app:linkdet}); nothing about the mask is specific to training.

\subsection{Text-attributed graphs}
\label{sec:method:tag}

A corpus is a directed graph $G=(V,E)$. Each node $v\in V$ is a document with a
token sequence $x_v$ and a string identifier (an article title, a
repository-qualified file path). An edge $(A,B)\in E$ exists when the text of $A$
contains a \emph{link} that resolves to $B$: a wikilink \texttt{[label](Title)}, an
\texttt{import} whose module path resolves to a file in the same repository, or a
\texttt{\textbackslash cite\{\}} whose key resolves to a paper in the corpus. Edges are
typed by modality and each modality has an \emph{online link detector} that, given a
token sequence, returns the set of links it contains as (\emph{end position},
\emph{target string}) pairs, where the end position is the token index one past the
link's closing delimiter. Detectors run at training time on packed sequences and at
generation time on the model's own output; they emit decoded strings and integer
offsets, so the mask machinery has no tokenizer dependency. Link \emph{positions} are
never stored in the graph---only that an edge exists---and are re-detected at runtime,
which forces the extractor that built the graph and the detector that reads it to
agree on what a link is (\Cref{app:datasets,app:linkdet}).

\subsection{Graph-traversal packing}
\label{sec:method:packing}

A training example is one packed sequence of exactly $T$ tokens ($T=32{,}768$
throughout) assembled from several documents. Each document occupies a contiguous
span laid out as an identifier card (a one-line comment carrying the document's
title or path, included with probability $\tfrac12$ per document and epoch so the
model sees both forms), the document body, and a single end-of-sequence token; only
bodies are ever truncated. Which documents share a pack is decided by a traversal of
$G$ from a uniformly sampled seed:
\begin{description}
  \item[BFS] expands the seed's out-neighbourhood level by level, in stored neighbour
        order (deterministic given the seed);
  \item[DFS] follows one out-path deep before backtracking, shuffling neighbours at each
        expansion;
  \item[random walk] a first-order walk over out-edges with a small uniform-teleport
        restart probability (not a restart-to-seed walk);
  \item[random] draws documents uniformly, ignoring $G$; the structure-free control.
\end{description}
A traversal is grown until the token budget is met, and a pack may hold several
disconnected neighbourhoods if the first is exhausted early. Because every seed is drawn
uniformly regardless of strategy, ``BFS packing'' means a set of uniformly seeded local
BFS neighbourhoods, not a global sweep (\Cref{app:traversal}).

\paragraph{Order matters.}
Attention is causal, so a document can only attend to documents packed \emph{before}
it. With out-edge traversal the natural proposal order puts a linking document ahead of
its target, which would leave every link unrealizable under the mask below. The packer
therefore reorders each connected component by a topological sort that places link
targets before the documents linking to them, and lays each component down as a
contiguous block. This targets-first reordering is applied identically under every
mask condition, including the controls, so it is not a confound in any contrast we
report; its consequence for what we \emph{can} measure is discussed in
\Cref{sec:analysis}.

\subsection{Link-aware attention masks}
\label{sec:method:masks}

Let $\mathrm{doc}(i)$ be the document occupying position $i\in[0,T)$ and
$\mathrm{comp}(i)$ the connected component of that document under the undirected
closure of $E$ restricted to the pack. A detected link $\ell$ in document $A$ with end
position $p_\ell$ and resolved target $B$ (present in the pack) defines a
\emph{grant}: the rectangle of (query, key) pairs
\[
  \mathcal{G}_\ell \;=\; [\,p_\ell,\ \mathrm{end}(A)\,)\times[\,\mathrm{start}(B),\ \mathrm{end}(B)\,),
\]
admitted only if $\mathrm{start}(B) < p_\ell$ (the target begins before the link---the
\emph{DAG gate}; a link to a later document is dropped, never granted acausally). The
four mask conditions, each a Boolean predicate $M(q,k)$ over query position $q$ and key
position $k$, are:
\begin{align}
  \doccausal:\quad  & M = [k\le q]\wedge[\mathrm{doc}(q)=\mathrm{doc}(k)] \\
  \crossdoc:\quad   & M = [k\le q]\wedge\bigl([\mathrm{doc}(q)=\mathrm{doc}(k)]\ \vee\ \textstyle\bigvee_{\ell}[(q,k)\in\mathcal{G}_\ell]\bigr) \\
  \concatlink:\quad & \text{as \crossdoc with } \mathcal{G}_\ell = [\,\mathrm{start}(A),\mathrm{end}(A)\,)\times[\,\mathrm{start}(B),\mathrm{end}(B)\,) \\
  \concat:\quad     & M = [k\le q]\wedge[\mathrm{comp}(q)=\mathrm{comp}(k)]
\end{align}
\doccausal is the baseline: block-diagonal causal attention, every document isolated,
equivalent to training on independent documents with the efficiency of packing.
\crossdoc is the method: when $A$ links to $B$, the tokens of $A$ \emph{from the link
onward} may read all of $B$. The grant is asymmetric ($B$ never sees $A$ unless it links
back), causality is never relaxed, and multiple grants compose by union. Because the
grant opens exactly where the link closes, the model learns that emitting a link is
what makes a document readable---the property the generation loop relies on.

The remaining two are \emph{compute controls}, ordered by the attention they permit
($\doccausal < \crossdoc \le \concatlink \le \concat$). \concatlink keeps the same
connectivity as \crossdoc but grants from the \emph{start} of the linking document
rather than from the link position: a strict superset of every \crossdoc grant, so it
isolates the effect of link-position gating from the effect of simply being able to
read the target. \concat merges each connected component into one causally attended
super-document: the most attention of the four, with no link supervision at all and
nothing an inference-time link could add. If \crossdoc beats both controls, the gain is
not explained by extra attention FLOPs or by mere co-location of linked content; that
residual---\crossdoc minus \concatlink at matched connectivity---is the number that
isolates the contribution (\Cref{sec:results}). \Cref{fig:masks} shows all four on one
packed batch.

\begin{figure}[t]
  \centering
  \begin{subfigure}[t]{0.48\linewidth}
    \includegraphics[width=\linewidth]{mask_doc_causal.png}
    \caption{\doccausal}
  \end{subfigure}\hfill
  \begin{subfigure}[t]{0.48\linewidth}
    \includegraphics[width=\linewidth]{mask_cross_doc_link.png}
    \caption{\crossdoc}
  \end{subfigure}\\[0.5em]
  \begin{subfigure}[t]{0.48\linewidth}
    \includegraphics[width=\linewidth]{mask_doc_concat_link.png}
    \caption{\concatlink}
  \end{subfigure}\hfill
  \begin{subfigure}[t]{0.48\linewidth}
    \includegraphics[width=\linewidth]{mask_doc_concatenated.png}
    \caption{\concat}
  \end{subfigure}
  \caption{The four attention masks on the same 16k-token pack (one Python repository
  from The Stack; dashed lines mark document boundaries, black is attendable).
  (a)~Every file attends only to itself. (b)~Where a file imports an earlier file, its
  rows extend left into the target---but only from the import line down, so each grant's
  top edge sits below the top of the importing file. (c)~The same grants extended to the
  top of the importing file (whole-document grant). (d)~All import-connected files merged
  into one causal triangle. Attention permitted increases from (a) to (d); (c) and (d)
  are compute controls for (b).}
  \label{fig:masks}
\end{figure}

\paragraph{Grant cap and block sparsity.}
Grants per pack are capped at $\texttt{max\_grants}=256$ in every reported run
(consumed in link-position order; later links beyond the cap are dropped with a
warning), and are encoded as bit-packed per-position masks so that membership is a
pointwise test with no $O(T^2)$ materialization. The mask is compiled to a block-sparse
form so fully masked blocks are skipped: PyTorch FlexAttention
\citep{flexattention_blog2024,flexattention_paper2024} serves as the reference path,
and training runs on custom Triton block-sparse kernels whose block schedule is computed
once per batch on the CPU and shared across all heads and layers. The kernels, the
grant encoding, and two numerical-correctness fixes they required are described in
\Cref{app:kernels}.

\subsection{Density-aware batch scheduling}
\label{sec:method:density}

Under \crossdoc the number of live attention blocks varies several-fold from pack to
pack---a pack of many short unlinked documents is nearly block-diagonal, one dominated
by a few long cross-linked documents is much denser---and the backward pass costs in
proportion. Under data-parallel training every rank must reach the gradient all-reduce
together, so each step runs at the speed of the densest pack in the group while the
other ranks idle. We remove this imbalance at the data-loading level. Packs are
precomputed offline for each epoch, with link detection done once in parallel workers;
each pack's density is scored analytically from its link positions by counting the
non-empty $128{\times}128$ blocks of its mask (exact under targets-first ordering, at
about a millisecond per pack), and packs are sorted into $32$ equal-count quantile
buckets. At each accumulation step every rank draws its pack from the \emph{same}
bucket---the bucket sequence is seeded by epoch, not by rank---so per-rank backward
cost is matched by construction, with no run-time cost measurement. The bucket
sequence shuffles across buckets so consecutive steps see different densities, and the
scheme bakes in no world size, so a run resumes at a different GPU count from a saved
schedule position. \Cref{fig:density-masks} shows the spread this equalizes; the
measured wall-clock effect is reported in \Cref{sec:results} and the mechanism and its
caveats in \Cref{app:density}.

\begin{figure}[t]
  \centering
  \includegraphics[width=\linewidth]{density_aware_masks.png}
  \caption{Block-level \crossdoc masks for one pack from the sparsest density bucket
  (left: many short documents, few links; almost all blocks empty) and one from the
  densest (right: a few large documents with cross-document grants filling large
  rectangles). Each cell is one $128$-token block, blue is non-empty, red lines mark
  document boundaries. Backward cost scales with the number of non-empty cells, so
  ranks holding packs from opposite ends of this range would otherwise wait on each
  other at every all-reduce.}
  \label{fig:density-masks}
\end{figure}

\subsection{Inference: link-following generation}
\label{sec:method:inference}

At inference the same detector runs over the model's own output after each generated
token (on a trailing window, firing when a link closes at the just-emitted token). A
detected target is resolved against the corpus by the \emph{same} key the training-time
matcher uses (exact identifier, then the detector's normalized key, then an optional
fuzzy title index). On a hit, the target document's tokens are inserted into the packed
context \emph{immediately before} the linking document---so that, exactly as in
training, the target starts before the link and the \crossdoc grant opens from the link
position into it---and generation continues with the mask rebuilt over the new layout.
A fetched document is itself scanned for links, giving bounded multi-hop retrieval to a
configurable depth (default two). Unresolved links may be skipped or, optionally,
satisfied by recursively \emph{generating} the missing document before the linking one,
which yields open-ended multi-document synthesis. This is retrieval by
\emph{insertion into the sequence} under the training mask, not by concatenating text
into a prompt: the fetched document is a first-class packed neighbour with the same
layout, span, and grant it would have had in a training pack. \Cref{alg:generation}
gives the loop. The current implementation recomputes the full forward pass at every
token (no key/value cache; inserting a document shifts every later rotary position,
which would invalidate a naive cache), a throughput limitation we discuss in
\Cref{sec:analysis}.

\begin{figure}[t]
\begin{center}
\fbox{\begin{minipage}{0.94\linewidth}\small
\textbf{Algorithm 1: link-following generation.}\\[2pt]
\textbf{Input:} corpus $C$, prompt for root document $r$, link detector $\mathcal{D}$,
max link depth $D$, retrieval mode.\\[2pt]
$\textsc{ctx} \gets [\,r\,]$ \hfill \textit{ordered list of document spans; root last}\\[2pt]
\textbf{procedure} \textsc{Generate}$(d,\ \mathit{depth})$:\\
\hspace*{1.5em}\textbf{while} $d$ is not finished:\\
\hspace*{3em}$S \gets \textsc{BuildSequence}(\textsc{ctx})$ \hfill \textit{concatenate spans; every target precedes its linker}\\
\hspace*{3em}$z \gets \textsc{Forward}(S,\ \text{\crossdoc mask})$ \hfill \textit{full forward; grants from detected links}\\
\hspace*{3em}$t \sim \textsc{Sample}(z \text{ at the last position of } d)$;\quad append $t$ to $d$\\
\hspace*{3em}\textbf{if} $\mathcal{D}$ fires on the trailing window, closing link $\ell$ with target string $s$, \textbf{and} $\mathit{depth} < D$:\\
\hspace*{4.5em}\textbf{if} $s$ resolves to $B \in C$: \hfill \textit{exact id $\to$ detector key $\to$ optional fuzzy index}\\
\hspace*{6em}insert $B$ into \textsc{ctx} immediately \emph{before} $d$ \hfill \textit{so the grant from $\ell$ opens into $B$}\\
\hspace*{6em}scan $B$ for its own links at $\mathit{depth}{+}1$ \hfill \textit{bounded multi-hop retrieval}\\
\hspace*{4.5em}\textbf{else if} generation fallback enabled:\\
\hspace*{6em}$B \gets$ new empty document before $d$;\quad \textsc{Generate}$(B,\ \mathit{depth}{+}1)$\\
\hspace*{3em}\textbf{if} $t = \texttt{EOS}$ or the per-document budget is exhausted: \textbf{stop}\\[2pt]
\textsc{Generate}$(r,\ 0)$
\end{minipage}}
\end{center}
\caption{Link-following generation. The detector, the resolution key, the
targets-before-linkers ordering, and the mask are the training-time implementations;
retrieval is insertion into the packed sequence, and an unresolved link may instead be
\emph{generated} recursively (open-ended multi-document synthesis). Context-window
pressure evicts the oldest non-active auxiliary documents.}
\label{alg:generation}
\end{figure}

\subsection{Model and training recipe}
\label{sec:method:recipe}

The backbone is a standard decoder-only transformer inherited from the open NanoGPT
speedrun lineage, held fixed across all mask conditions so that any measured
difference is attributable to the attention structure. Every component of it is prior
art; \Cref{app:recipe} documents the recipe and \Cref{sec:experiments} the tuned
operating point. One choice matters to the method itself: rotary positions run over
the global packed index with no per-document reset (\Cref{sec:analysis}).
